
\begin{remark}[A weight-doubling witness sampler at the free-coordinate grades; AI-checked, pending human review]
\label{rem:fc-witness-h3}
This remark comes from an AI-only pass (\file{research/audits/reviews/2026-11-05-rees-s3-witness-h3.md}, with code in its folder), checked and corrected by a second, adversarial AI pass with independent code (\file{research/audits/reviews/2026-11-05-rees-s3-witness-h3-check.md}); no human has checked it. (R1), Theorem~\ref{thm:free-coordinate} and the tables are unchanged, and only the sentence after Theorem~\ref{thm:free-coordinate} on the classical time of part~(iii) with these nodes uses it. No novelty is claimed: the sampler is the construction of Remark~\ref{rem:fc-zero-sum-arity}(a) with another digit set. Let $d \in \Gamma_{\mathrm F}$, write $\kappa = \kappa_d$ and $h = (p-1)/\kappa \ge 2$, and let the labels be independent and uniform on $K_d^\perp$, as in Lemma~\ref{lem:fc-forest}(a). The witness entries lie in $D = P_\kappa\cup\{0\}$, the $h$-th roots of unity together with $0$.

\textup{(a)} \emph{Depth (proved; AI-checked, pending human review).} Put $D_0 = D$ and $D_{j+1} = D_j - D_j$, and let $j^*_h$ be the least $J$ with $D_J = \Fp$; at $h = 1$ this is the $j^*$ of Remark~\ref{rem:fc-zero-sum-arity}. Then $0, 1 \in D_j$ and $|D_j| \ge \min(p, 2^jh+1)$ by the Cauchy--Davenport theorem, so $j^*_h \le \lceil\log_2\kappa\rceil$. If $\sqrt p > (\kappa-1)(\kappa-2)+3$, then $D_0 - D_0 = \Fp$ and $j^*_h = 1$: for $y \ne 0$, Jacobi sums show that at least $(p - (\kappa-1)(\kappa-2)\sqrt p - 3\kappa + 1)/\kappa^2$ pairs $(u,v) \in P_\kappa^2$ have $u - v = y$, and if $\sqrt p \ge 2(\kappa-1)(\kappa-2)+4$, a uniform $u \in P_\kappa$ has $u - y \in P_\kappa$ with probability at least $1/(2\kappa)$. For fixed $h$, let $\varphi(h)$ be Euler's function and $c_h$ the largest $\ell_1$-norm of a power of $\zeta_h$ in the power basis of $\mathbb{Z}[\zeta_h]$. Then $\log_2((p^{1/\varphi(h)}-1)/(2c_h)) \le j^*_h \le \lceil\log_2(c_h\varphi(h)^2p^{1/\varphi(h)}/2)\rceil + 1$: every element of $D_j$ is the image of a signed sum of at most $2^j$ powers of $\zeta_h$, and a Minkowski vector of the prime $(p,\zeta_h-\omega)$, $\omega \in \Fp$ of order $h$, gives short representatives of every residue. So $j^*_h = \log_2p/\varphi(h) + O_h(1)$. Computed values (two independent implementations): $j^*_3 = 1,1,2,2$ at $p = 7,13,19,31$, $j^*_4 = 1,2$ at $p = 13,17$, and $j^*_3 - \tfrac12\log_2p \in [-0.85, 0.60]$ for the primes $p \equiv 1 \pmod 3$ below $6\cdot 10^4$.

\textup{(b)} \emph{Sampler and node (proved; AI-checked, pending human review).} The level-$j$ solvers of Remark~\ref{rem:fc-zero-sum-arity}(a) work with $D_j$ in place of $D^{(j)}$, with $j^*_h$ in place of $j^*$, and with a fixed table writing each $\alpha \in D_{j+1}$ as $x - x'$ with $x,x' \in D_j$ (and $0 = 0 - 0$). The proof is unchanged, since it uses only $1 \in D_j$, the table, the nonzero kernel vector at level $j^*_h$, and the uniformity of each $\beta_k$ given everything read before. The level-$0$ solver returns nonzero weights in $P_\kappa$ on $R$ columns with $\mathbb{E}R = N_0(n)$, where $N_j$ is defined as there; $N_0$ is a polynomial of exact degree $1 + j^*_h$ in $n$, with leading coefficient $\sigma_p^{\,j^*_h}/(1+j^*_h)!$. With running sums the time is $O(R(n^2+p) + j^*_hp^2)$ operations in $\Fp$, the last term for the tables. When $\sqrt p \ge 2(\kappa-1)(\kappa-2)+4$ the solver is used in a rescaled form that needs no table. In dimension $n' \ge 1$, for $k = 0,\dots,\lceil\log_2p\rceil-1$, it draws $n'$ fresh columns until the kernel vector $\alpha$ of their first $n'-1$ coordinates has $\beta = \sum_i\alpha_iv_{i,n'} \ne 0$, and rescales $\alpha$ to $2^k\alpha/\beta$. It splits each entry $e$ as $x - x'$ with $x = b^\kappa$ for a uniform $b \in \Fp^\times$ and $x' = x - e \in D$, with at most $\lceil 2\kappa\ln 2\,(B + \lceil\log_2M\rceil)\rceil$ trials per entry, and lets the binary digits of the residue replace the covering loop. Its expected number of columns is exactly $1 + \lceil\log_2p\rceil\frac{p}{p-1}\frac{n(n+1)}{2}$, and all splits succeed except with probability $2^{-B}$ when at most $M$ entries are split. The node at $d$ is that of Remark~\ref{rem:fc-zero-sum-arity}(b) with this solver, followed by (F4)--(F6). It uses the rescaled form when $\sqrt p \ge 2(\kappa-1)(\kappa-2)+4$, with its exact column count in place of $\bar N_0$. Otherwise it uses the tables; then $p < (2(\kappa-1)(\kappa-2)+4)^2$, so $j^*_h$ and the tables cost time polynomial in $\kappa$. It never returns a wrong witness, and it aborts with probability at most $\tau + 3\cdot2^{-B}$, the last term for the capped root extractions, field draws and splitting searches. Its arity is $M_d = O((2\log_2p+5)^{j^*_h}g_d^{\,1+j^*_h}\log(1/\tau))$, of degree $1 + j^*_h \le 1 + \lceil\log_2\kappa_d\rceil$ in $g_d$, and of degree $2$ when $\sqrt p > (\kappa_d-1)(\kappa_d-2)+3$. Its time is polynomial in $M_d$, $\kappa_d$ and $\log p$. In finite computations (AI-run, pending human review) at $h = 3$ ($p = 7,13,19,31$) and $h = 4$ ($p = 13,17$), two independent implementations checked every output exactly, their mean column counts matched $N_0$, and the node aborted in at most $1/40$ of trials at $\tau = 1/20$, never with a wrong witness; the rescaled form ran at $p = 2^{31}-1$ with $\kappa = 3$ and at $p = 2^{61}-1$ with $\kappa = 6$. For comparison at fixed $\kappa$: Schedule~1 has degree $\kappa\lceil\log_2(\kappa+1)\rceil$ ($12$ against $2$ at $p = 13$, $\kappa = 4$, a prime of $E_2(4)$), and Theorem~\ref{thm:power-sde}(ii) has degree $2$ but needs $p \ge 16\kappa^{2\kappa+2}$; both are worst-case and use no label law.

\textup{(c)} \emph{Fixed class (proof of the sentence after Theorem~\ref{thm:free-coordinate}; AI-checked, pending human review).} Under the hypotheses of part~(iii), use the node of~(b) at every grade of $\Gamma_{\mathrm F}$, the node of Remark~\ref{rem:fc-zero-sum-arity}(b) at $\Gamma_0$ and the linear node at $\Gamma_1$, with $\tau_d = \eps/(9(c-1)Sn_d)$ at the grades of $\Gamma_{\mathrm F}\cup\Gamma_0$. Each such node reads only its own bank, checks exactly and converts by Proposition~\ref{prop:descent}, so it is a further case of Lemma~\ref{lem:fc-forest}(a), like the zero-sum node. It contributes at most $\tau_d + 3\cdot2^{-B}$ to Lemma~\ref{lem:fc-forest}(b), so the bound $2\eps/9$ there and the rest of the proof of part~(iii) are unchanged: the recursion returns $\hh$ with probability at least $1-\eps$ from $c\,S\prod_dM_d$ queries. Since $n_d = \prod_{d'<d}M_{d'}$, $\log(1/\tau_d) = O_c(\log(m/\eps) + \log\log p)$. At fixed $c$, $\kappa_d \le c$, so every $M_d$, the number of nodes and the time of each node are polynomial in $m$, $\log p$ and $\log(1/\eps)$, for every prime $p > c$. Count $g_d \le m$ and ignore factors polylogarithmic in $m$, $p$ and $1/\eps$. Then the exponent of $m$ in the query count is $E = 1 + \sum_{d=2}^c(1 + j_d)$, with $j_d = j^*_{h_d}$ at $\Gamma_{\mathrm F}$, $j_d = \lceil\log_2(p-1)\rceil$ at $\Gamma_0$ and $j_d = 0$ at $\Gamma_1$. For $p > \max_{2\le\kappa\le c}((\kappa-1)(\kappa-2)+3)^2$, every $j_d$ at $\Gamma_{\mathrm F}$ is $1$ and $\Gamma_0$ is empty, so $E \le 2c-1$. Scanning the smaller primes gives $\max_{p>c}E = 3, 6, 8, 10, 13, 14, 16, 19$ for $c = 2,\dots,9$ (computed by two independent implementations; at $c = 9$ the maximum is attained at $p = 19$). At the class-nine profile of Table~\ref{tab:fcs-budget}, with the arities of~(b) as stated and not optimized, $\log_2Q$ lies between $93.0$ and $150.8$ at the six primes of that table, against $34.0$ to $42.3$ for part~(iii) (computed by one implementation, AI-run): these nodes trade queries for classical time. The exponent $E$ concerns a procedure with aborts on the i.i.d.\ supply only. It is not compared here with $E_W(c)$ of Theorem~\ref{thm:bb-envelope}, which concerns worst-case rules that use no label law.

\textup{(d)} \emph{Limits (proved; AI-checked, pending human review).} At arity $M$ (labels read from the independent uniform supply), a union bound over the nonzero $y \in D^M$ shows that a witness exists with probability at most $((h+1)^M-1)/p^{a_d} < (h+1)^{M-a_d}q^{a_d}$. So no method of any running time succeeds with probability $1/2$ below $a_d\log_{h+1}p - \log_{h+1}2$ labels, and at arity $a_d + O(1)$ a witness exists with probability bounded below only if $h \ge \kappa^{\Omega(a_d)}$; the arity of Lemma~\ref{lem:fc-choice} is of the order of that threshold. Consider a switch hierarchy whose level-$(j+1)$ relations, multiplied by one scalar per level, are differences of level-$j$ weightings, and whose base accepts every kernel vector. By induction its level-$j$ digit set lies in $\gamma D_j$ for some $\gamma \in \Fp^\times$, so its depth is at least $j^*_h$. For fixed $h \ge 3$ its depth, and the degree $1 + j^*_h$ of the sampler in~(b), therefore grow like $\log_2p/\varphi(h)$. Hierarchies with a separate scalar for each switch are not covered, and whether a deep hierarchy can have bounded degree is not settled here. On a uniform bank the witness problem is Module-SIS over $\mathbb{Z}[\zeta_h]$ modulo the degree-one prime $(p,\zeta_h-\omega)$, with solutions in $\mu_h\cup\{0\}$. For even $h$ these are, by Kronecker's theorem, exactly the vectors of canonical $\ell_\infty$-norm at most $1$, and at $h = 2$ this is the ternary SIS of Remark~\ref{rem:fc-zero-sum-sis}. No worst-case hardness result for this modulus is used or claimed. Whether some method has degree bounded uniformly in $p$ for fixed $h \ge 3$ remains open.
% NEEDS HUMAN VERIFICATION: Jacobi-sum absolute values (Ireland--Rosen, ch. 8), Kronecker's theorem on algebraic integers with all conjugates of modulus at most 1, and the absence of a worst-case reduction for Module-SIS modulo a degree-one prime ideal (Langlois--Stehle 2015 treat the modulus qR), all from memory; no bib entries are added. sigma_p <= ceil(2 log2(p-1)) + 4 is the bound imported in rem:fc-zero-sum-arity(a); with sigma_p <= p instead (proved there) only the constants of the tables case change, since there p < (2(kappa-1)(kappa-2)+4)^2.
\end{remark}
